% APA Format for Universidad de La Salle
\documentclass[a4paper,12pt]{report}
\usepackage[T1]{fontenc}
\usepackage[{{language}}]{babel}
\usepackage{acronym}
\usepackage{amsfonts}
\usepackage{amsmath}
\usepackage{amssymb}
\usepackage{array}
\usepackage{booktabs}
\usepackage{caption}
\usepackage{enumitem}
\usepackage{fancyhdr}
\usepackage{float}
\usepackage{geometry}
\usepackage{graphicx}
\usepackage{hyperref}
\usepackage{indentfirst}
\usepackage{lmodern}
\usepackage{longtable}
\usepackage{makecell}
\usepackage{microtype}
\usepackage[nottoc]{tocbibind}
\usepackage{siunitx}
\usepackage{titlesec}
\usepackage{url}
\usepackage{csquotes}

\geometry{top=3cm,bottom=3cm,left=3cm,right=3cm}
\setlength{\emergencystretch}{2em}

\titleformat{\chapter}[display]{\normalfont\huge\bfseries}{\chaptertitlename\ \thechapter}{20pt}{\Huge}
\titlespacing*{\chapter}{0pt}{-30pt}{20pt}
\titleformat{\section}[block]{\normalfont\Large\bfseries}{\thesection}{1em}{}
\titleformat{\subsection}[block]{\normalfont\large\bfseries}{\thesubsection}{1em}{}

% APA tables and figures: the number is bold and the title is italic on its
% own line, both placed above the object; a table note sits below the table.
\captionsetup{position=top,labelsep=newline,labelfont=bf,textfont=it,
  justification=raggedright,singlelinecheck=false,skip=4pt}
\AtBeginDocument{\iflanguage{spanish}{\captionsetup[table]{name=Tabla}}{}}

\numberwithin{equation}{chapter}
% APA numbers tables and figures through the document, not by chapter.
\counterwithout{figure}{chapter}
\counterwithout{table}{chapter}

\pagestyle{fancy}
\fancyhf{}
\cfoot{\thepage}

\newcommand{\titulo}{{{title}}}
\newcommand{\autor}{{{author}}}
\newcommand{\codigo}{{{code}}}
\newcommand{\programa}{{{program}}}
\newcommand{\universidad}{Universidad de La Salle}
\newcommand{\director}{{{director}}}
\newcommand{\correo}{{{email}}}

\hypersetup{pdfauthor={\autor},pdftitle={\titulo},pdfsubject={\programa},
  pdfkeywords={investigación},colorlinks=true,linkcolor=black,citecolor=black,urlcolor=black}

\begin{document}

  \begin{titlepage}
    \centering\vspace*{\fill}
    \includegraphics[width=8cm]{assets/images/logosalle.png}\\[1cm]
    \universidad\\[0.5cm]\programa\\[2cm]
    {\LARGE\bfseries \titulo}\\[2cm]
    \textbf{Research Proposal}\\[2cm]
    \Large\autor\\\codigo\\[1cm]
    Advisor: \director\\\correo\\[2cm]
    Submission Date: \today\\\vspace*{\fill}
  \end{titlepage}

  \thispagestyle{empty}\newpage

  % Replace this abstract with a one-paragraph summary of your own proposal.
  \chapter*{Abstract}
  Write your abstract here. This document provides a template for academic projects in APA format.
  \textit{Keywords}: keyword1, keyword2, keyword3

  \chapter{Introduction}\label{chap:introduction}
  % Replace this paragraph with the background and the problem of your own study.
  Provide background context and describe the research problem; the sample
  proposal builds on the framing given in~\cite{example}.

  \section{Objectives}
  % Replace these sample objectives with your own.
  \begin{enumerate}
    \item Objective 1
    \item Objective 2
    \item Objective 3
  \end{enumerate}

  \chapter{Methodology}\label{chap:methodology}
  % Replace this paragraph with your own design, population and instruments.
  Describe the research approach, the population and sample, and the data
  collection methods. The standard procedure behind the sample study is
  presented in a textbook~\cite{example2}. The diagram below draws the
  general flow of the methodology; replace it with a diagram of your own.

  % Diagram drawn from its source when the document is built, so the project
  % carries no image file. Keep style= and the caption; wrap a caption that
  % contains a comma in braces.
  \begin{mermaid}[style=monochrome, caption={General flow of the methodology}, pos=H]
    flowchart LR
    plan[Plan the study] --> collect[Collect data] --> analyze[Analyze data]
  \end{mermaid}

  \chapter{Expected Results}\label{chap:results}
  % Replace this paragraph and the table with your own expectations and data.
  Describe the expected results and how they will be summarized; an online
  note illustrates the same kind of report~\cite{example3}.

  % APA table: the bold number and the italic title go above the table and the
  % note goes below it. Replace the rows, the caption and the label.
  \begin{table}[H]
    \caption{Illustrative expected values}
    \label{tab:expected}
    \begin{tabular}{@{}lcc@{}}
      \toprule
      Case & Initial & Final \\
      \midrule
      A & 12 & 8 \\
      B & 20 & 11 \\
      C & 15 & 9 \\
      \bottomrule
    \end{tabular}
    \par\medskip
    \begin{minipage}{\linewidth}
      \footnotesize\textit{Note.} Example values; replace them with your own data.
    \end{minipage}
  \end{table}

  \chapter{Conclusions and Future Work}\label{chap:conclusions}
  % Replace this paragraph with your own summary and next steps.
  Summarize the contribution of the proposal, name its limitations and state
  what should be done next.

  % APA reference list. Every entry in bib/references.bib is cited above.
  \renewcommand{\bibname}{References}
  \bibliographystyle{apalike}
  \bibliography{bib/references}

  \appendix
  \chapter{Example code}
  % Replace this listing with the code that produced the table above.
  Keep it short: at most twelve lines, and set lang= to the language of the
  block.

  \begin{code}[lang=python, numbers=true]
def mean(values):
    total = sum(values)
    count = len(values)
    return total / count if count else 0
  \end{code}

\end{document}
